% ====================================================================
% РАЗДЕЛ: МЕХАНИКА (КИНЕМАТИКА)
% ФАЙЛ: tasks/01_mechanics/kim01.tex
% ИСТОЧНИК: Демидова ЕГЭ 2026, Вариант 1–20, Задание №1
% ====================================================================

% === ЗАДАЧА 1 ===
% Тег: #МЕХАНИКА #КИМ_01 #ВАРИАНТ_01
\begin{taskbasic}[code={\label{task:dem26_v1_n1}}]
\textbf{Задача 1.} На рисунке приведён график зависимости от времени $t$ проекции $v_x$ скорости тела, движущегося прямолинейно вдоль оси $Ox$[cite: 1]. Определите проекцию $s_x$ перемещения этого тела в интервале времени от $0$ до $15$~с[cite: 1]. Ответ запишите с учётом знака проекции[cite: 1].

\nopagebreak\smallskip
\begin{center}
\begin{tikzpicture}[x=0.35cm, y=0.08cm, every node/.style={font=\footnotesize}]
    % Сетка (шаг 5 по x, шаг 10 по y)
    \draw[xstep=5, ystep=10, gray!30, thin] (0,-20) grid (15,20);
    % Оси
    \draw[-{Stealth[scale=0.8]}, black, thick] (0,0) -- (17,0) node[right] {$t, \text{ с}$};
    \draw[-{Stealth[scale=0.8]}, black, thick] (0,-20) -- (0,24) node[above] {$v_x, \text{ м/с}$};
    % Оцифровка
    \node[left] at (0,20) {$20$};
    \node[left] at (0,10) {$10$};
    \node[below left] at (0,0) {$0$};
    \node[left] at (0,-10) {$-10$};
    \node[left] at (0,-20) {$-20$};
    \node[below] at (5,0) {$5$};
    \node[below] at (10,0) {$10$};
    \node[below] at (15,0) {$15$};
    % График
    \draw[ultra thick, brandPrimary] (0,-20) -- (10,0) -- (15,10);
\end{tikzpicture}
\end{center}
\kim{1}
\end{taskbasic}
\answer{-75~м}

% === ЗАДАЧА 2 ===
% Тег: #МЕХАНИКА #КИМ_01 #ВАРИАНТ_02
\begin{taskbasic}[code={\label{task:dem26_v2_n1}}]
\textbf{Задача 2.} На рисунке приведён график зависимости от времени $t$ проекции $v_x$ скорости тела, движущегося прямолинейно вдоль оси $Ox$[cite: 1]. Определите проекцию $s_x$ перемещения этого тела в интервале времени от $5$ до $15$~с[cite: 1]. Ответ запишите с учётом знака проекции[cite: 1].

\nopagebreak\smallskip
\begin{center}
\begin{tikzpicture}[x=0.35cm, y=0.08cm, every node/.style={font=\footnotesize}]
    % Сетка
    \draw[xstep=5, ystep=10, gray!30, thin] (0,-20) grid (15,20);
    % Оси
    \draw[-{Stealth[scale=0.8]}, black, thick] (0,0) -- (17,0) node[right] {$t, \text{ с}$};
    \draw[-{Stealth[scale=0.8]}, black, thick] (0,-20) -- (0,24) node[above] {$v_x, \text{ м/с}$};
    % Оцифровка
    \node[left] at (0,20) {$20$};
    \node[left] at (0,10) {$10$};
    \node[below left] at (0,0) {$0$};
    \node[left] at (0,-10) {$-10$};
    \node[left] at (0,-20) {$-20$};
    \node[below] at (5,0) {$5$};
    \node[below] at (10,0) {$10$};
    \node[below] at (15,0) {$15$};
    % График
    \draw[ultra thick, brandPrimary] (0,-20) -- (10,0) -- (15,10);
\end{tikzpicture}
\end{center}
\kim{1}
\end{taskbasic}
\answer{-25~м}

% === ЗАДАЧА 3 ===
% Тег: #МЕХАНИКА #КИМ_01 #ВАРИАНТ_03
\begin{taskbasic}[code={\label{task:dem26_v3_n1}}]
\textbf{Задача 3.} На рисунке приведён график зависимости проекции $v_x$ скорости тела от времени $t$[cite: 1]. Определите проекцию $a_x$ ускорения тела в интервале времени от $20$ до $30$~с[cite: 1]. Ответ запишите с учётом знака проекции[cite: 1].

\nopagebreak\smallskip
\begin{center}
\begin{tikzpicture}[x=0.15cm, y=0.12cm, every node/.style={font=\footnotesize}]
    % Сетка (шаг 10 по x, шаг 10 по y)
    \draw[xstep=10, ystep=10, gray!30, thin] (0,0) grid (40,25);
    % Оси
    \draw[-{Stealth[scale=0.8]}, black, thick] (0,0) -- (44,0) node[right] {$t, \text{ с}$};
    \draw[-{Stealth[scale=0.8]}, black, thick] (0,0) -- (0,27) node[above] {$v_x, \text{ м/с}$};
    % Оцифровка
    \node[left] at (0,20) {$20$};
    \node[left] at (0,10) {$10$};
    \node[below left] at (0,0) {$0$};
    \node[below] at (10,0) {$10$};
    \node[below] at (20,0) {$20$};
    \node[below] at (30,0) {$30$};
    \node[below] at (40,0) {$40$};
    % График
    \draw[ultra thick, brandPrimary] (0,5) -- (10,20) -- (20,10) -- (30,20) -- (40,0);
\end{tikzpicture}
\end{center}
\kim{1}
\end{taskbasic}
\answer{1~м/с$^{2}$}

% === ЗАДАЧА 4 ===
% Тег: #МЕХАНИКА #КИМ_01 #ВАРИАНТ_04
\begin{taskbasic}[code={\label{task:dem26_v4_n1}}]
\textbf{Задача 4.} На рисунке приведён график зависимости проекции $v_x$ скорости тела от времени $t$[cite: 1]. Определите проекцию $a_x$ ускорения тела в интервале времени от $30$ до $40$~с[cite: 1]. Ответ запишите с учётом знака проекции[cite: 1].

\nopagebreak\smallskip
\begin{center}
\begin{tikzpicture}[x=0.15cm, y=0.12cm, every node/.style={font=\footnotesize}]
    % Сетка
    \draw[xstep=10, ystep=10, gray!30, thin] (0,0) grid (40,30);
    % Оси
    \draw[-{Stealth[scale=0.8]}, black, thick] (0,0) -- (44,0) node[right] {$t, \text{ с}$};
    \draw[-{Stealth[scale=0.8]}, black, thick] (0,0) -- (0,32) node[above] {$v_x, \text{ м/с}$};
    % Оцифровка
    \node[left] at (0,30) {$30$};
    \node[left] at (0,20) {$20$};
    \node[left] at (0,10) {$10$};
    \node[below left] at (0,0) {$0$};
    \node[below] at (10,0) {$10$};
    \node[below] at (20,0) {$20$};
    \node[below] at (30,0) {$30$};
    \node[below] at (40,0) {$40$};
    % График
    \draw[ultra thick, brandPrimary] (0,5) -- (10,20) -- (20,10) -- (30,25) -- (40,0);
\end{tikzpicture}
\end{center}
\kim{1}
\end{taskbasic}
\answer{-2{,}5~м/с$^{2}$}

% === ЗАДАЧА 5 ===
% Тег: #МЕХАНИКА #КИМ_01 #ВАРИАНТ_05
\begin{taskbasic}[code={\label{task:dem26_v5_n1}}]
\textbf{Задача 5.} Тело движется вдоль оси $Ox$[cite: 1]. На рисунке приведён график зависимости проекции $v_x$ скорости тела от времени $t$[cite: 1]. Определите среднюю скорость тела в интервале времени от $20$ до $40$~с[cite: 1].

\nopagebreak\smallskip
\begin{center}
\begin{tikzpicture}[x=0.15cm, y=0.12cm, every node/.style={font=\footnotesize}]
    % Сетка
    \draw[xstep=10, ystep=10, gray!30, thin] (0,0) grid (40,30);
    % Оси
    \draw[-{Stealth[scale=0.8]}, black, thick] (0,0) -- (44,0) node[right] {$t, \text{ с}$};
    \draw[-{Stealth[scale=0.8]}, black, thick] (0,0) -- (0,35) node[above] {$v_x, \text{ м/с}$};
    % Оцифровка
    \node[left] at (0,30) {$30$};
    \node[left] at (0,20) {$20$};
    \node[left] at (0,10) {$10$};
    \node[below left] at (0,0) {$0$};
    \node[below] at (10,0) {$10$};
    \node[below] at (20,0) {$20$};
    \node[below] at (30,0) {$30$};
    \node[below] at (40,0) {$40$};
    % График
    \draw[ultra thick, brandPrimary] (0,15) -- (10,5) -- (20,20) -- (30,0) -- (40,25);
\end{tikzpicture}
\end{center}
\kim{1}
\end{taskbasic}
\answer{10~м/с}

% === ЗАДАЧА 6 ===
% Тег: #МЕХАНИКА #КИМ_01 #ВАРИАНТ_06
\begin{taskbasic}[code={\label{task:dem26_v6_n1}}]
\textbf{Задача 6.} Тело движется вдоль оси $Ox$[cite: 1]. На рисунке приведён график зависимости проекции $v_x$ скорости тела от времени $t$[cite: 1]. Определите среднюю скорость тела в интервале времени от $0$ до $20$~с[cite: 1].

\nopagebreak\smallskip
\begin{center}
\begin{tikzpicture}[x=0.15cm, y=0.12cm, every node/.style={font=\footnotesize}]
    % Сетка
    \draw[xstep=10, ystep=10, gray!30, thin] (0,0) grid (40,25);
    % Оси
    \draw[-{Stealth[scale=0.8]}, black, thick] (0,0) -- (44,0) node[right] {$t, \text{ с}$};
    \draw[-{Stealth[scale=0.8]}, black, thick] (0,0) -- (0,27) node[above] {$v_x, \text{ м/с}$};
    % Оцифровка
    \node[left] at (0,20) {$20$};
    \node[left] at (0,10) {$10$};
    \node[below left] at (0,0) {$0$};
    \node[below] at (10,0) {$10$};
    \node[below] at (20,0) {$20$};
    \node[below] at (30,0) {$30$};
    \node[below] at (40,0) {$40$};
    % График
    \draw[ultra thick, brandPrimary] (0,15) -- (10,5) -- (20,20) -- (30,0) -- (40,25);
\end{tikzpicture}
\end{center}
\kim{1}
\end{taskbasic}
\answer{11{,}25~м/с}

% === ЗАДАЧА 7 ===
% Тег: #МЕХАНИКА #КИМ_01 #ВАРИАНТ_07
\begin{taskbasic}[code={\label{task:dem26_v7_n1}}]
\textbf{Задача 7.} На рисунке приведён график зависимости проекции $v_x$ скорости тела от времени $t$[cite: 1]. Определите путь, пройденный телом, в интервале времени от $20$ до $40$~с[cite: 1].

\nopagebreak\smallskip
\begin{center}
\begin{tikzpicture}[x=0.15cm, y=0.12cm, every node/.style={font=\footnotesize}]
    % Сетка (шаг 10 по x, шаг 10 по y)
    \draw[xstep=10, ystep=10, gray!30, thin] (0,0) grid (40,25);
    % Оси
    \draw[-{Stealth[scale=0.8]}, black, thick] (0,0) -- (44,0) node[right] {$t, \text{ с}$};
    \draw[-{Stealth[scale=0.8]}, black, thick] (0,0) -- (0,27) node[above] {$v_x, \text{ м/с}$};
    % Оцифровка
    \node[left] at (0,20) {$20$};
    \node[left] at (0,10) {$10$};
    \node[below left] at (0,0) {$0$};
    \node[below] at (10,0) {$10$};
    \node[below] at (20,0) {$20$};
    \node[below] at (30,0) {$30$};
    \node[below] at (40,0) {$40$};
    % График
    \draw[ultra thick, brandPrimary] (0,5) -- (10,20) -- (20,10) -- (30,25) -- (40,0);
\end{tikzpicture}
\end{center}
\kim{1}
\end{taskbasic}
\answer{250~м}

% === ЗАДАЧА 8 ===
% Тег: #МЕХАНИКА #КИМ_01 #ВАРИАНТ_08
\begin{taskbasic}[code={\label{task:dem26_v8_n1}}]
\textbf{Задача 8.} На рисунке приведён график зависимости проекции $v_x$ скорости тела от времени $t$[cite: 1]. Определите путь, пройденный телом в интервале времени от $0$ до $20$~с[cite: 1].

\nopagebreak\smallskip
\begin{center}
\begin{tikzpicture}[x=0.15cm, y=0.12cm, every node/.style={font=\footnotesize}]
    % Сетка
    \draw[xstep=10, ystep=10, gray!30, thin] (0,0) grid (40,25);
    % Оси
    \draw[-{Stealth[scale=0.8]}, black, thick] (0,0) -- (44,0) node[right] {$t, \text{ с}$};
    \draw[-{Stealth[scale=0.8]}, black, thick] (0,0) -- (0,27) node[above] {$v_x, \text{ м/с}$};
    % Оцифровка
    \node[left] at (0,20) {$20$};
    \node[left] at (0,10) {$10$};
    \node[below left] at (0,0) {$0$};
    \node[below] at (10,0) {$10$};
    \node[below] at (20,0) {$20$};
    \node[below] at (30,0) {$30$};
    \node[below] at (40,0) {$40$};
    % График
    \draw[ultra thick, brandPrimary] (0,5) -- (10,20) -- (20,10) -- (30,25) -- (40,0);
\end{tikzpicture}
\end{center}
\kim{1}
\end{taskbasic}
\answer{275~м}

% === ЗАДАЧА 9 ===
% Тег: #МЕХАНИКА #КИМ_01 #ВАРИАНТ_09
\begin{taskbasic}[code={\label{task:dem26_v9_n1}}]
\textbf{Задача 9.} На рисунке приведён график зависимости проекции $v_x$ скорости тела от времени $t$[cite: 1]. Определите путь, пройденный телом в течение интервала времени от $0$ до $3$~с[cite: 1].

\nopagebreak\smallskip
\begin{center}
\begin{tikzpicture}[x=0.35cm, y=0.08cm, every node/.style={font=\footnotesize}]
    % Сетка (шаг 2 по x, шаг 10 по y)
    \draw[xstep=2, ystep=10, gray!30, thin] (0,-10) grid (11,20);
    % Оси
    \draw[-{Stealth[scale=0.8]}, black, thick] (0,0) -- (12,0) node[right] {$t, \text{ с}$};
    \draw[-{Stealth[scale=0.8]}, black, thick] (0,-10) -- (0,24) node[above] {$v_x, \text{ м/с}$};
    % Оцифровка
    \node[left] at (0,20) {$20$};
    \node[left] at (0,10) {$10$};
    \node[below left] at (0,0) {$0$};
    \node[left] at (0,-10) {$-10$};
    \node[below] at (2,0) {$2$};
    \node[below] at (4,0) {$4$};
    \node[below] at (6,0) {$6$};
    \node[below] at (8,0) {$8$};
    \node[below] at (10,0) {$10$};
    % График
    \draw[ultra thick, brandPrimary] (0,20) -- (3,-10) -- (7,0) -- (11,10);
\end{tikzpicture}
\end{center}
\kim{1}
\end{taskbasic}
\answer{35~м}

% === ЗАДАЧА 10 ===
% Тег: #МЕХАНИКА #КИМ_01 #ВАРИАНТ_10
\begin{taskbasic}[code={\label{task:dem26_v10_n1}}]
\textbf{Задача 10.} На рисунке приведён график зависимости проекции $v_x$ скорости тела от времени $t$. Определите проекцию $a_x$ ускорения этого тела в интервале времени от $3$ до $7$~с.

\nopagebreak\smallskip
\begin{center}
\begin{tikzpicture}[x=0.35cm, y=0.08cm, every node/.style={font=\footnotesize}]
    % Сетка
    \draw[xstep=2, ystep=10, gray!30, thin] (0,-10) grid (11,20);
    % Оси
    \draw[-{Stealth[scale=0.8]}, black, thick] (0,0) -- (12,0) node[right] {$t, \text{ с}$};
    \draw[-{Stealth[scale=0.8]}, black, thick] (0,-10) -- (0,24) node[above] {$v_x, \text{ м/с}$};
    % Оцифровка
    \node[left] at (0,20) {$20$};
    \node[left] at (0,10) {$10$};
    \node[below left] at (0,0) {$0$};
    \node[left] at (0,-10) {$-10$};
    \node[below] at (2,0) {$2$};
    \node[below] at (4,0) {$4$};
    \node[below] at (6,0) {$6$};
    \node[below] at (8,0) {$8$};
    \node[below] at (10,0) {$10$};
    % График
    \draw[ultra thick, brandPrimary] (0,20) -- (3,-10) -- (7,0) -- (11,10);
\end{tikzpicture}
\end{center}
\kim{1}
\end{taskbasic}
\answer{2{,}5~м/с$^{2}$}

% === ЗАДАЧА 11 ===
% Тег: #МЕХАНИКА #КИМ_01 #ВАРИАНТ_11
\begin{taskbasic}[code={\label{task:dem26_v11_n1}}]
\textbf{Задача 11.} На рисунке приведён график зависимости проекции $v_x$ скорости тела от времени $t$[cite: 5]. Определите проекцию $a_x$ ускорения этого тела в момент времени $25$~с[cite: 5]. Ответ запишите с учётом знака проекции[cite: 5].

\nopagebreak\smallskip
\begin{center}
\begin{tikzpicture}[x=0.25cm, y=0.12cm, every node/.style={font=\footnotesize}]
    % Сетка
    \draw[xstep=5, ystep=5, gray!30, thin] (0,-15) grid (35,15);
    % Оси
    \draw[-{Stealth[scale=0.8]}, black, thick] (0,0) -- (38,0) node[right] {$t, \text{ с}$};
    \draw[-{Stealth[scale=0.8]}, black, thick] (0,-15) -- (0,18) node[above] {$v_x, \text{ м/с}$};
    % Оцифровка
    \node[left] at (0,15) {$15$};
    \node[left] at (0,10) {$10$};
    \node[left] at (0,5) {$5$};
    \node[below left] at (0,0) {$0$};
    \node[left] at (0,-5) {$-5$};
    \node[left] at (0,-10) {$-10$};
    \node[left] at (0,-15) {$-15$};
    \node[below] at (10,0) {$10$};
    \node[below] at (20,0) {$20$};
    \node[below] at (30,0) {$30$};
    % График
    \draw[ultra thick, brandPrimary] (0,0) -- (10,-15) -- (20,-15) -- (35,15);
\end{tikzpicture}
\end{center}
\kim{1}
\end{taskbasic}
\answer{2~м/с$^{2}$}

% === ЗАДАЧА 12 ===
% Тег: #МЕХАНИКА #КИМ_01 #ВАРИАНТ_12
\begin{taskbasic}[code={\label{task:dem26_v12_n1}}]
\textbf{Задача 12.} На рисунке приведён график зависимости проекции $v_x$ скорости тела от времени $t$. Определите проекцию $s_x$ перемещения этого тела в интервале времени от $0$ до $35$~с. Ответ запишите с учётом знака проекции.

\nopagebreak\smallskip
\begin{center}
\begin{tikzpicture}[x=0.25cm, y=0.12cm, every node/.style={font=\footnotesize}]
    % Сетка
    \draw[xstep=5, ystep=5, gray!30, thin] (0,-15) grid (35,15);
    % Оси
    \draw[-{Stealth[scale=0.8]}, black, thick] (0,0) -- (38,0) node[right] {$t, \text{ с}$};
    \draw[-{Stealth[scale=0.8]}, black, thick] (0,-15) -- (0,18) node[above] {$v_x, \text{ м/с}$};
    % Оцифровка
    \node[left] at (0,15) {$15$};
    \node[left] at (0,10) {$10$};
    \node[left] at (0,5) {$5$};
    \node[below left] at (0,0) {$0$};
    \node[left] at (0,-5) {$-5$};
    \node[left] at (0,-10) {$-10$};
    \node[left] at (0,-15) {$-15$};
    \node[below] at (10,0) {$10$};
    \node[below] at (20,0) {$20$};
    \node[below] at (30,0) {$30$};
    % График
    \draw[ultra thick, brandPrimary] (0,0) -- (10,-15) -- (20,-15) -- (35,15);
\end{tikzpicture}
\end{center}
\kim{1}
\end{taskbasic}
\answer{-187{,}5~м}

% === ЗАДАЧА 13 ===
% Тег: #МЕХАНИКА #КИМ_01 #ВАРИАНТ_13
\begin{taskbasic}[code={\label{task:dem26_v13_n1}}]
\textbf{Задача 13.} На рисунке показан график зависимости проекции $v_x$ скорости тела от времени $t$. Какова проекция $a_x$ ускорения этого тела в момент времени $t = 4{,}5$~с? Ответ запишите с учётом знака проекции.

\nopagebreak\smallskip
\begin{center}
\begin{tikzpicture}[x=0.8cm, y=0.25cm, every node/.style={font=\footnotesize}]
    % Сетка
    \draw[xstep=1, ystep=4, gray!30, thin] (0,-12) grid (6,12);
    % Оси
    \draw[-{Stealth[scale=0.8]}, black, thick] (0,0) -- (6.5,0) node[right] {$t, \text{ с}$};
    \draw[-{Stealth[scale=0.8]}, black, thick] (0,-12) -- (0,14) node[above] {$v_x, \text{ м/с}$};
    % Оцифровка
    \node[left] at (0,12) {$12$};
    \node[left] at (0,8) {$8$};
    \node[left] at (0,4) {$4$};
    \node[below left] at (0,0) {$0$};
    \node[left] at (0,-4) {$-4$};
    \node[left] at (0,-8) {$-8$};
    \node[left] at (0,-12) {$-12$};
    \node[below] at (2,0) {$2$};
    \node[below] at (4,0) {$4$};
    \node[below] at (6,0) {$6$};
    % График
    \draw[ultra thick, brandPrimary] (0,-12) -- (2,-4) -- (3,12) -- (4,12) -- (5,0) -- (6,8);
\end{tikzpicture}
\end{center}
\kim{1}
\end{taskbasic}
\answer{-12~м/с$^{2}$}

% === ЗАДАЧА 14 ===
% Тег: #МЕХАНИКА #КИМ_01 #ВАРИАНТ_14
\begin{taskbasic}[code={\label{task:dem26_v14_n1}}]
\textbf{Задача 14.} На рисунке показан график зависимости проекции $v_x$ скорости тела от времени $t$. Какова проекция $a_x$ ускорения этого тела в момент времени $t = 1$~с? Ответ запишите с учётом знака проекции.

\nopagebreak\smallskip
\begin{center}
\begin{tikzpicture}[x=0.8cm, y=0.25cm, every node/.style={font=\footnotesize}]
    % Сетка
    \draw[xstep=1, ystep=4, gray!30, thin] (0,-12) grid (6,12);
    % Оси
    \draw[-{Stealth[scale=0.8]}, black, thick] (0,0) -- (6.5,0) node[right] {$t, \text{ с}$};
    \draw[-{Stealth[scale=0.8]}, black, thick] (0,-12) -- (0,14) node[above] {$v_x, \text{ м/с}$};
    % Оцифровка
    \node[left] at (0,12) {$12$};
    \node[left] at (0,8) {$8$};
    \node[left] at (0,4) {$4$};
    \node[below left] at (0,0) {$0$};
    \node[left] at (0,-4) {$-4$};
    \node[left] at (0,-8) {$-8$};
    \node[left] at (0,-12) {$-12$};
    \node[below] at (2,0) {$2$};
    \node[below] at (4,0) {$4$};
    \node[below] at (6,0) {$6$};
    % График
    \draw[ultra thick, brandPrimary] (0,-12) -- (2,-4) -- (3,12) -- (4,12) -- (5,0) -- (6,8);
\end{tikzpicture}
\end{center}
\kim{1}
\end{taskbasic}
\answer{4~м/с$^{2}$}

% === ЗАДАЧА 15 ===
% Тег: #МЕХАНИКА #КИМ_01 #ВАРИАНТ_15
\begin{taskbasic}[code={\label{task:dem26_v15_n1}}]
\textbf{Задача 15.} На рисунке представлен график зависимости модуля скорости $v$ тела от времени $t$. Найдите среднюю скорость тела за время от $0$ до $10$~с.

\nopagebreak\smallskip
\begin{center}
\begin{tikzpicture}[x=0.45cm, y=0.2cm, every node/.style={font=\footnotesize}]
    % Сетка
    \draw[xstep=2, ystep=5, gray!30, thin] (0,0) grid (12,15);
    % Оси
    \draw[-{Stealth[scale=0.8]}, black, thick] (0,0) -- (13,0) node[right] {$t, \text{ с}$};
    \draw[-{Stealth[scale=0.8]}, black, thick] (0,0) -- (0,17) node[above] {$v, \text{ м/с}$};
    % Оцифровка
    \node[left] at (0,15) {$15$};
    \node[left] at (0,10) {$10$};
    \node[left] at (0,5) {$5$};
    \node[below left] at (0,0) {$0$};
    \node[below] at (2,0) {$2$};
    \node[below] at (4,0) {$4$};
    \node[below] at (6,0) {$6$};
    \node[below] at (8,0) {$8$};
    \node[below] at (10,0) {$10$};
    \node[below] at (12,0) {$12$};
    % График
    \draw[ultra thick, brandPrimary] (0,0) -- (4,15) -- (12,15);
\end{tikzpicture}
\end{center}
\kim{1}
\end{taskbasic}
\answer{12~м/с}

% === ЗАДАЧА 16 ===
% Тег: #МЕХАНИКА #КИМ_01 #ВАРИАНТ_16
\begin{taskbasic}[code={\label{task:dem26_v16_n1}}]
\textbf{Задача 16.} На рисунке представлен график зависимости модуля скорости $v$ тела от времени $t$. Найдите путь, пройденный телом за время от $0$ до $12$~с.

\nopagebreak\smallskip
\begin{center}
\begin{tikzpicture}[x=0.45cm, y=0.2cm, every node/.style={font=\footnotesize}]
    % Сетка
    \draw[xstep=2, ystep=5, gray!30, thin] (0,0) grid (12,15);
    % Оси
    \draw[-{Stealth[scale=0.8]}, black, thick] (0,0) -- (13,0) node[right] {$t, \text{ с}$};
    \draw[-{Stealth[scale=0.8]}, black, thick] (0,0) -- (0,17) node[above] {$v, \text{ м/с}$};
    % Оцифровка
    \node[left] at (0,15) {$15$};
    \node[left] at (0,10) {$10$};
    \node[left] at (0,5) {$5$};
    \node[below left] at (0,0) {$0$};
    \node[below] at (2,0) {$2$};
    \node[below] at (4,0) {$4$};
    \node[below] at (6,0) {$6$};
    \node[below] at (8,0) {$8$};
    \node[below] at (10,0) {$10$};
    \node[below] at (12,0) {$12$};
    % График
    \draw[ultra thick, brandPrimary] (0,0) -- (4,15) -- (12,15);
\end{tikzpicture}
\end{center}
\kim{1}
\end{taskbasic}
\answer{150~м}

% === ЗАДАЧА 17 ===
% Тег: #МЕХАНИКА #КИМ_01 #ВАРИАНТ_17
% Из графика варианта 17 в присланном PDF: отрезок от (0,12) до (2,4)
\begin{taskbasic}[code={\label{task:dem26_v17_n1}}]
\textbf{Задача 17.} На рисунке показан график зависимости проекции $v_x$ скорости тела от времени $t$. Какова проекция $a_x$ ускорения этого тела в момент времени $4{,}5$~с? Ответ запишите с учётом знака проекции.

\nopagebreak\smallskip
\begin{center}
\begin{tikzpicture}[x=0.8cm, y=0.25cm, every node/.style={font=\footnotesize}]
    % Сетка
    \draw[xstep=1, ystep=4, gray!30, thin] (0,-12) grid (6,12);
    % Оси
    \draw[-{Stealth[scale=0.8]}, black, thick] (0,0) -- (6.5,0) node[right] {$t, \text{ с}$};
    \draw[-{Stealth[scale=0.8]}, black, thick] (0,-12) -- (0,14) node[above] {$v_x, \text{ м/с}$};
    % Оцифровка
    \node[left] at (0,12) {$12$};
    \node[left] at (0,8) {$8$};
    \node[left] at (0,4) {$4$};
    \node[below left] at (0,0) {$0$};
    \node[left] at (0,-4) {$-4$};
    \node[left] at (0,-8) {$-8$};
    \node[left] at (0,-12) {$-12$};
    \node[below] at (2,0) {$2$};
    \node[below] at (4,0) {$4$};
    \node[below] at (6,0) {$6$};
    % График
    \draw[ultra thick, brandPrimary] (0,12) -- (2,4) -- (3,-12) -- (4,-12) -- (5,0) -- (6,-8);
\end{tikzpicture}
\end{center}
\kim{1}
\end{taskbasic}
\answer{12~м/с$^{2}$}

% === ЗАДАЧА 18 ===
% Тег: #МЕХАНИКА #КИМ_01 #ВАРИАНТ_18
\begin{taskbasic}[code={\label{task:dem26_v18_n1}}]
\textbf{Задача 18.} На рисунке показан график зависимости проекции $v_x$ скорости тела от времени $t$. Определите путь, пройденный этим телом в интервале времени от $3$ до $6$~с.

\nopagebreak\smallskip
\begin{center}
\begin{tikzpicture}[x=0.8cm, y=0.25cm, every node/.style={font=\footnotesize}]
    % Сетка
    \draw[xstep=1, ystep=4, gray!30, thin] (0,-12) grid (6,12);
    % Оси
    \draw[-{Stealth[scale=0.8]}, black, thick] (0,0) -- (6.5,0) node[right] {$t, \text{ с}$};
    \draw[-{Stealth[scale=0.8]}, black, thick] (0,-12) -- (0,14) node[above] {$v_x, \text{ м/с}$};
    % Оцифровка
    \node[left] at (0,12) {$12$};
    \node[left] at (0,8) {$8$};
    \node[left] at (0,4) {$4$};
    \node[below left] at (0,0) {$0$};
    \node[left] at (0,-4) {$-4$};
    \node[left] at (0,-8) {$-8$};
    \node[left] at (0,-12) {$-12$};
    \node[below] at (2,0) {$2$};
    \node[below] at (4,0) {$4$};
    \node[below] at (6,0) {$6$};
    % График
    \draw[ultra thick, brandPrimary] (0,12) -- (2,4) -- (3,-12) -- (4,-12) -- (5,0) -- (6,-8);
\end{tikzpicture}
\end{center}
\kim{1}
\end{taskbasic}
\answer{16~м}

% === ЗАДАЧА 19 ===
% Тег: #МЕХАНИКА #КИМ_01 #ВАРИАНТ_19
\begin{taskbasic}[code={\label{task:dem26_v19_n1}}]
\textbf{Задача 19.} На рисунке показан график зависимости проекции $v_x$ скорости тела от времени $t$. Какова проекция $s_x$ перемещения этого тела за время от момента $t_1 = 0{,}5$~с до момента $t_2 = 1{,}5$~с? Ответ запишите с учётом знака проекции.

\nopagebreak\smallskip
\begin{center}
\begin{tikzpicture}[x=1.6cm, y=0.25cm, every node/.style={font=\footnotesize}]
    % Сетка
    \draw[xstep=0.5, ystep=2, gray!30, thin] (0,-6) grid (3,6);
    % Оси
    \draw[-{Stealth[scale=0.8]}, black, thick] (0,0) -- (3.3,0) node[right] {$t, \text{ с}$};
    \draw[-{Stealth[scale=0.8]}, black, thick] (0,-6) -- (0,7) node[above] {$v_x, \text{ м/с}$};
    % Оцифровка
    \node[left] at (0,6) {$6$};
    \node[left] at (0,4) {$4$};
    \node[left] at (0,2) {$2$};
    \node[below left] at (0,0) {$0$};
    \node[left] at (0,-2) {$-2$};
    \node[left] at (0,-4) {$-4$};
    \node[left] at (0,-6) {$-6$};
    \node[below] at (1,0) {$1$};
    \node[below] at (2,0) {$2$};
    \node[below] at (3,0) {$3$};
    % График
    \draw[ultra thick, brandPrimary] (0,4) -- (0.5,0) -- (1,6) -- (1.5,2) -- (2.5,2) -- (3,-6);
\end{tikzpicture}
\end{center}
\kim{1}
\end{taskbasic}
\answer{3~м}

% === ЗАДАЧА 20 ===
% Тег: #МЕХАНИКА #КИМ_01 #ВАРИАНТ_20
\begin{taskbasic}[code={\label{task:dem26_v20_n1}}]
\textbf{Задача 20.} На рисунке показан график зависимости проекции $v_x$ скорости тела от времени $t$. Какова проекция $a_x$ ускорения этого тела в момент времени $2$~с? Ответ запишите с учётом знака проекции.

\nopagebreak\smallskip
\begin{center}
\begin{tikzpicture}[x=1.6cm, y=0.25cm, every node/.style={font=\footnotesize}]
    % Сетка
    \draw[xstep=0.5, ystep=2, gray!30, thin] (0,-6) grid (3,6);
    % Оси
    \draw[-{Stealth[scale=0.8]}, black, thick] (0,0) -- (3.3,0) node[right] {$t, \text{ с}$};
    \draw[-{Stealth[scale=0.8]}, black, thick] (0,-6) -- (0,7) node[above] {$v_x, \text{ м/с}$};
    % Оцифровка
    \node[left] at (0,6) {$6$};
    \node[left] at (0,4) {$4$};
    \node[left] at (0,2) {$2$};
    \node[below left] at (0,0) {$0$};
    \node[left] at (0,-2) {$-2$};
    \node[left] at (0,-4) {$-4$};
    \node[left] at (0,-6) {$-6$};
    \node[below] at (1,0) {$1$};
    \node[below] at (2,0) {$2$};
    \node[below] at (3,0) {$3$};
    % График
    \draw[ultra thick, brandPrimary] (0,4) -- (0.5,0) -- (1,6) -- (1.5,2) -- (2.5,2) -- (3,-6);
\end{tikzpicture}
\end{center}
\kim{1}
\end{taskbasic}
\answer{0~м/с$^{2}$}

% === ЗАДАЧА 21 ===
% Тег: #МЕХАНИКА #КИМ_01 #ВАРИАНТ_21
\begin{taskbasic}[code={\label{task:dem26_v21_n1}}]
\textbf{Задача 21.} На рисунке приведён график зависимости от времени $t$ проекции $v_x$ скорости тела, движущегося прямолинейно вдоль оси $Ox$. Определите проекцию $a_x$ ускорения этого тела в интервале времени от $15$ до $20$~с. Ответ запишите с учётом знака проекции.
\begin{center}
\begin{tikzpicture}[x=0.35cm, y=0.12cm, every node/.style={font=\footnotesize}]
    % Сетка
    \draw[xstep=5, ystep=5, gray!30, thin] (0,-20) grid (30,20);
    % Оси
    \draw[-{Stealth[scale=0.8]}, black, thick] (0,0) -- (33,0) node[right] {$t, \text{ с}$};
    \draw[-{Stealth[scale=0.8]}, black, thick] (0,-25) -- (0,23) node[above] {$v_x, \text{ м/с}$};
    % Оцифровка
    \node[left] at (0,20) {$20$};
    \node[left] at (0,10) {$10$};
    \node[below left] at (0,0) {$0$};
    \node[left] at (0,-10) {$-10$};
    \node[left] at (0,-20) {$-20$};
    \node[below] at (5,0) {$5$};
    \node[below] at (10,0) {$10$};
    \node[below] at (15,0) {$15$};
    \node[below] at (20,0) {$20$};
    \node[below] at (25,0) {$25$};
    \node[below] at (30,0) {$30$};
    % График
    \draw[ultra thick, brandPrimary] (0,-20) -- (10,0) -- (15,10) -- (20,-10) -- (30,20);
\end{tikzpicture}
\end{center}
\kim{1}
\end{taskbasic}
\answer{-4~м/с$^{2}$}

% === ЗАДАЧА 22 ===
% Тег: #МЕХАНИКА #КИМ_01 #ВАРИАНТ_22
\begin{taskbasic}[code={\label{task:dem26_v22_n1}}]
\textbf{Задача 22.} На рисунке приведён график зависимости от времени $t$ проекции $v_x$ тела, движущегося прямолинейно вдоль оси $Ox$. Определите проекцию $a_x$ ускорения этого тела в интервале времени от $20$ до $30$~с. Ответ запишите с учётом знака проекции.
\begin{center}
\begin{tikzpicture}[x=0.35cm, y=0.12cm, every node/.style={font=\footnotesize}]
    % Сетка
    \draw[xstep=5, ystep=5, gray!30, thin] (0,-20) grid (30,20);
    % Оси
    \draw[-{Stealth[scale=0.8]}, black, thick] (0,0) -- (33,0) node[right] {$t, \text{ с}$};
    \draw[-{Stealth[scale=0.8]}, black, thick] (0,-25) -- (0,23) node[above] {$v_x, \text{ м/с}$};
    % Оцифровка
    \node[left] at (0,20) {$20$};
    \node[left] at (0,10) {$10$};
    \node[below left] at (0,0) {$0$};
    \node[left] at (0,-10) {$-10$};
    \node[left] at (0,-20) {$-20$};
    \node[below] at (5,0) {$5$};
    \node[below] at (10,0) {$10$};
    \node[below] at (15,0) {$15$};
    \node[below] at (20,0) {$20$};
    \node[below] at (25,0) {$25$};
    \node[below] at (30,0) {$30$};
    % График
    \draw[ultra thick, brandPrimary] (0,-20) -- (10,0) -- (15,10) -- (20,-10) -- (30,20);
\end{tikzpicture}
\end{center}
\kim{1}
\end{taskbasic}
\answer{3~м/с$^{2}$}

% === ЗАДАЧА 23 ===
% Тег: #МЕХАНИКА #КИМ_01 #ВАРИАНТ_23
\begin{taskbasic}[code={\label{task:dem26_v23_n1}}]
\textbf{Задача 23.} На рисунке приведён график зависимости проекции $v_x$ скорости тела от времени $t$. Определите проекцию $a_x$ ускорения этого тела в интервале времени от $0$ до $5$~с. Ответ запишите с учётом знака проекции.
\begin{center}
\begin{tikzpicture}[x=0.35cm, y=0.15cm, every node/.style={font=\footnotesize}]
    % Сетка
    \draw[xstep=5, ystep=5, gray!30, thin] (0,-20) grid (30,15);
    % Оси
    \draw[-{Stealth[scale=0.8]}, black, thick] (0,0) -- (33,0) node[right] {$t, \text{ с}$};
    \draw[-{Stealth[scale=0.8]}, black, thick] (0,-20) -- (0,18) node[above] {$v_x, \text{ м/с}$};
    % Оцифровка
    \node[left] at (0,15) {$15$};
    \node[left] at (0,10) {$10$};
    \node[left] at (0,5) {$5$};
    \node[below left] at (0,0) {$0$};
    \node[left] at (0,-5) {$-5$};
    \node[left] at (0,-10) {$-10$};
    \node[left] at (0,-15) {$-15$};
    \node[left] at (0,-20) {$-20$};
    \node[below] at (5,0) {$5$};
    \node[below] at (10,0) {$10$};
    \node[below] at (15,0) {$15$};
    \node[below] at (20,0) {$20$};
    \node[below] at (25,0) {$25$};
    \node[below] at (30,0) {$30$};
    % График
    \draw[ultra thick, brandPrimary] (0,15) -- (5,0) -- (10,-15) -- (15,10) -- (20,0) -- (30,-20);
\end{tikzpicture}
\end{center}
\kim{1}
\end{taskbasic}
\answer{-3~м/с$^{2}$}

% === ЗАДАЧА 24 ===
% Тег: #МЕХАНИКА #КИМ_01 #ВАРИАНТ_24
\begin{taskbasic}[code={\label{task:dem26_v24_n1}}]
\textbf{Задача 24.} На рисунке приведён график зависимости проекции $v_x$ скорости тела от времени $t$. Определите проекцию $a_x$ ускорения этого тела в интервале времени от $20$ до $30$~с. Ответ запишите с учётом знака проекции.
\begin{center}
\begin{tikzpicture}[x=0.35cm, y=0.15cm, every node/.style={font=\footnotesize}]
    % Сетка
    \draw[xstep=5, ystep=5, gray!30, thin] (0,-20) grid (30,15);
    % Оси
    \draw[-{Stealth[scale=0.8]}, black, thick] (0,0) -- (33,0) node[right] {$t, \text{ с}$};
    \draw[-{Stealth[scale=0.8]}, black, thick] (0,-20) -- (0,18) node[above] {$v_x, \text{ м/с}$};
    % Оцифровка
    \node[left] at (0,15) {$15$};
    \node[left] at (0,10) {$10$};
    \node[left] at (0,5) {$5$};
    \node[below left] at (0,0) {$0$};
    \node[left] at (0,-5) {$-5$};
    \node[left] at (0,-10) {$-10$};
    \node[left] at (0,-15) {$-15$};
    \node[left] at (0,-20) {$-20$};
    \node[below] at (5,0) {$5$};
    \node[below] at (10,0) {$10$};
    \node[below] at (15,0) {$15$};
    \node[below] at (20,0) {$20$};
    \node[below] at (25,0) {$25$};
    \node[below] at (30,0) {$30$};
    % График
    \draw[ultra thick, brandPrimary] (0,15) -- (5,0) -- (10,-15) -- (15,10) -- (20,0) -- (30,-20);
\end{tikzpicture}
\end{center}
\kim{1}
\end{taskbasic}
\answer{-2~м/с$^{2}$}

% === ЗАДАЧА 25 ===
% Тег: #МЕХАНИКА #КИМ_01 #ВАРИАНТ_25
\begin{taskbasic}[code={\label{task:dem26_v25_n1}}]
\textbf{Задача 25.} На рисунке приведён график зависимости координаты тела $x$ от времени $t$ при прямолинейном движении тела вдоль оси $Ox$. Определите проекцию перемещения этого тела на ось $Ox$ в промежутке времени от $0$ до $14$~с. Ответ запишите с учётом знака проекции.
\begin{center}
\begin{tikzpicture}[x=0.45cm, y=0.15cm, every node/.style={font=\footnotesize}]
    % Сетка
    \draw[xstep=2, ystep=5, gray!30, thin] (0,-15) grid (14,15);
    % Оси
    \draw[-{Stealth[scale=0.8]}, black, thick] (0,0) -- (15.5,0) node[right] {$t, \text{ с}$};
    \draw[-{Stealth[scale=0.8]}, black, thick] (0,-15) -- (0,18) node[above] {$x, \text{ м}$};
    % Оцифровка
    \node[left] at (0,15) {$15$};
    \node[left] at (0,10) {$10$};
    \node[left] at (0,5) {$5$};
    \node[below left] at (0,0) {$0$};
    \node[left] at (0,-5) {$-5$};
    \node[left] at (0,-10) {$-10$};
    \node[left] at (0,-15) {$-15$};
    \foreach \x in {2,4,6,8,10,12,14} \node[below] at (\x,0) {\x};
    % График
    \draw[ultra thick, brandPrimary] (0,-10) -- (4,-10) -- (8,10) -- (14,-5);
\end{tikzpicture}
\end{center}
\kim{1}
\end{taskbasic}
\answer{5~м}

% === ЗАДАЧА 26 ===
% Тег: #МЕХАНИКА #КИМ_01 #ВАРИАНТ_26
\begin{taskbasic}[code={\label{task:dem26_v26_n1}}]
\textbf{Задача 26.} На рисунке приведён график зависимости координаты тела $x$ от времени $t$ при прямолинейном движении тела вдоль оси $Ox$. Определите проекцию скорости этого тела на ось $Ox$ в промежутке времени от $8$ до $12$~с. Ответ запишите с учётом знака проекции.
\begin{center}
\begin{tikzpicture}[x=0.45cm, y=0.15cm, every node/.style={font=\footnotesize}]
    % Сетка
    \draw[xstep=2, ystep=5, gray!30, thin] (0,-15) grid (14,15);
    % Оси
    \draw[-{Stealth[scale=0.8]}, black, thick] (0,0) -- (15.5,0) node[right] {$t, \text{ с}$};
    \draw[-{Stealth[scale=0.8]}, black, thick] (0,-15) -- (0,18) node[above] {$x, \text{ м}$};
    % Оцифровка
    \node[left] at (0,15) {$15$};
    \node[left] at (0,10) {$10$};
    \node[left] at (0,5) {$5$};
    \node[below left] at (0,0) {$0$};
    \node[left] at (0,-5) {$-5$};
    \node[left] at (0,-10) {$-10$};
    \node[left] at (0,-15) {$-15$};
    \foreach \x in {2,4,6,8,10,12,14} \node[below] at (\x,0) {\x};
    % График
    \draw[ultra thick, brandPrimary] (0,-10) -- (4,-10) -- (8,10) -- (14,-5);
\end{tikzpicture}
\end{center}
\kim{1}
\end{taskbasic}
\answer{-2,5~м/с}

% === ЗАДАЧА 27 ===
% Тег: #МЕХАНИКА #КИМ_01 #ВАРИАНТ_27
\begin{taskbasic}[code={\label{task:dem26_v27_n1}}]
\textbf{Задача 27.} Тело движется вдоль оси $Ox$. На рисунке приведён график зависимости проекции $v_x$ скорости тела от времени $t$. Определите путь, пройденный телом в интервале времени от $12$ до $18$~с.
\begin{center}
\begin{tikzpicture}[x=0.35cm, y=0.15cm, every node/.style={font=\footnotesize}]
    % Сетка
    \draw[xstep=5, ystep=5, gray!30, thin] (0,-10) grid (20,15);
    % Оси
    \draw[-{Stealth[scale=0.8]}, black, thick] (0,0) -- (22,0) node[right] {$t, \text{ с}$};
    \draw[-{Stealth[scale=0.8]}, black, thick] (0,-10) -- (0,18) node[above] {$v_x, \text{ м/с}$};
    % Оцифровка
    \node[left] at (0,15) {$15$};
    \node[left] at (0,10) {$10$};
    \node[left] at (0,5) {$5$};
    \node[below left] at (0,0) {$0$};
    \node[left] at (0,-5) {$-5$};
    \node[left] at (0,-10) {$-10$};
    \node[below] at (5,0) {$5$};
    \node[below] at (10,0) {$10$};
    \node[below] at (15,0) {$15$};
    \node[below] at (20,0) {$20$};
    % График
    \draw[ultra thick, brandPrimary] (0,5) -- (3,5) -- (5,15) -- (10,-10) -- (13,-10) -- (18,5) -- (20,5);
\end{tikzpicture}
\end{center}
\kim{1}
\end{taskbasic}
\answer{25~м}

% === ЗАДАЧА 28 ===
% Тег: #МЕХАНИКА #КИМ_01 #ВАРИАНТ_28
\begin{taskbasic}[code={\label{task:dem26_v28_n1}}]
\textbf{Задача 28.} Тело движется вдоль оси $Ox$. На рисунке приведён график зависимости проекции $v_x$ скорости тела от времени $t$. Определите проекцию перемещения тела $s_x$ в интервале времени от $5$ до $10$~с. Ответ запишите с учётом знака проекции.
\begin{center}
\begin{tikzpicture}[x=0.35cm, y=0.15cm, every node/.style={font=\footnotesize}]
    % Сетка
    \draw[xstep=5, ystep=5, gray!30, thin] (0,-10) grid (20,15);
    % Оси
    \draw[-{Stealth[scale=0.8]}, black, thick] (0,0) -- (22,0) node[right] {$t, \text{ с}$};
    \draw[-{Stealth[scale=0.8]}, black, thick] (0,-10) -- (0,18) node[above] {$v_x, \text{ м/с}$};
    % Оцифровка
    \node[left] at (0,15) {$15$};
    \node[left] at (0,10) {$10$};
    \node[left] at (0,5) {$5$};
    \node[below left] at (0,0) {$0$};
    \node[left] at (0,-5) {$-5$};
    \node[left] at (0,-10) {$-10$};
    \node[below] at (5,0) {$5$};
    \node[below] at (10,0) {$10$};
    \node[below] at (15,0) {$15$};
    \node[below] at (20,0) {$20$};
    % График
    \draw[ultra thick, brandPrimary] (0,5) -- (3,5) -- (5,15) -- (10,-10) -- (13,-10) -- (18,5) -- (20,5);
\end{tikzpicture}
\end{center}
\kim{1}
\end{taskbasic}
\answer{12{,}5~м}

% === ЗАДАЧА 29 ===
% Тег: #МЕХАНИКА #КИМ_01 #ВАРИАНТ_29
\begin{taskbasic}[code={\label{task:dem26_v29_n1}}]
\textbf{Задача 29.} На рисунке приведён график зависимости проекции $v_x$ скорости тела от времени $t$. Определите проекцию $a_x$ ускорения этого тела в интервале времени от $0$ до $10$~с. Ответ запишите с учётом знака проекции.
\begin{center}
\begin{tikzpicture}[x=0.35cm, y=0.12cm, every node/.style={font=\footnotesize}]
    % Сетка
    \draw[xstep=5, ystep=5, gray!30, thin] (0,-20) grid (30,20);
    % Оси
    \draw[-{Stealth[scale=0.8]}, black, thick] (0,0) -- (33,0) node[right] {$t, \text{ с}$};
    \draw[-{Stealth[scale=0.8]}, black, thick] (0,-20) -- (0,24) node[above] {$v_x, \text{ м/с}$};
    % Оцифровка
    \node[left] at (0,20) {$20$};
    \node[left] at (0,10) {$10$};
    \node[below left] at (0,0) {$0$};
    \node[left] at (0,-10) {$-10$};
    \node[left] at (0,-20) {$-20$};
    \node[below] at (5,0) {$5$};
    \node[below] at (10,0) {$10$};
    \node[below] at (15,0) {$15$};
    \node[below] at (20,0) {$20$};
    \node[below] at (25,0) {$25$};
    \node[below] at (30,0) {$30$};
    % График
    \draw[ultra thick, brandPrimary] (0,-20) -- (10,0) -- (15,10) -- (20,-10) -- (30,20);
\end{tikzpicture}
\end{center}
\kim{1}
\end{taskbasic}
\answer{2~м/с$^{2}$}

% === ЗАДАЧА 30 ===
% Тег: #МЕХАНИКА #КИМ_01 #ВАРИАНТ_30
\begin{taskbasic}[code={\label{task:dem26_v30_n1}}]
\textbf{Задача 30.} На рисунке приведён график зависимости проекции $v_x$ скорости тела от времени $t$. Определите проекцию $a_x$ ускорения этого тела в интервале времени от $15$ до $20$~с. Ответ запишите с учётом знака проекции.
\begin{center}
\begin{tikzpicture}[x=0.35cm, y=0.12cm, every node/.style={font=\footnotesize}]
    % Сетка
    \draw[xstep=5, ystep=5, gray!30, thin] (0,-20) grid (30,20);
    % Оси
    \draw[-{Stealth[scale=0.8]}, black, thick] (0,0) -- (33,0) node[right] {$t, \text{ с}$};
    \draw[-{Stealth[scale=0.8]}, black, thick] (0,-20) -- (0,24) node[above] {$v_x, \text{ м/с}$};
    % Оцифровка
    \node[left] at (0,20) {$20$};
    \node[left] at (0,10) {$10$};
    \node[below left] at (0,0) {$0$};
    \node[left] at (0,-10) {$-10$};
    \node[left] at (0,-20) {$-20$};
    \node[below] at (5,0) {$5$};
    \node[below] at (10,0) {$10$};
    \node[below] at (15,0) {$15$};
    \node[below] at (20,0) {$20$};
    \node[below] at (25,0) {$25$};
    \node[below] at (30,0) {$30$};
    % График
    \draw[ultra thick, brandPrimary] (0,-20) -- (10,0) -- (15,10) -- (20,-10) -- (30,20);
\end{tikzpicture}
\end{center}
\kim{1}
\end{taskbasic}
\answer{-4~м/с$^{2}$}